理想 Michelson 干涉仪的两个臂产生可调光程差 Δ（已经包含双程因素），由中心角频率 ω₀>0 的准单色平稳光源照明。光源功率谱关于 ω₀ 呈窄带高斯分布 S(ω)∝exp[−(ω−ω₀)²/(2σ²)]，σ>0 且 σ≪ω₀；忽略负频尾、色散和偏振失配，两臂返回光强相等。
(1) 对单一角频率写出相位差与输出干涉强度；(2) 对高斯光谱积分求总强度和条纹可见度；(3) 定义可见度下降为 1/e 时的光程差为相干长度，求它与 σ 的关系；(4) 讨论改变 Δ 如何从干涉数据恢复光源谱宽以及 σ→0 极限。
